% ---------------------------------------------------------------------------
% DRAFT 8 — Judge methodology (design, described before results land) and the
% significance/tiers section update (36/15 pairs + clustered-bootstrap robustness).
% [[JUDGE_RESULTS_PENDING]] marks every number that depends on the still-running
% phi4-mini full-coverage pass. Do not submit with this marker present.
% ---------------------------------------------------------------------------

\subsection{Cross-Model Judge}
\label{sec:judge}

An LLM-as-judge audit is only as convincing as the judge's independence from what it is judging.
The item-level judge used to compute judge-audited accuracy is Gemini 2.5 Flash---one of the twelve
evaluated models. To assess whether this affects the audit, we additionally judge every REG/NUM/TMP
prediction with \textbf{phi4-mini} (Microsoft Phi, 3.8B parameters), run locally via Ollama at zero
cost. Phi shares no model family with any of the twelve benchmarked systems. We call this a
\textbf{cross-model} judge, not an independent or ground-truth one: phi4-mini is a small model, and
disagreement between it and Gemini 2.5 Flash is compatible with either judge being wrong, which is
why we pair it with human adjudication rather than treating its verdict as final.

The cross-model pass covers \textbf{every} REG/NUM/TMP prediction---344 items $\times$ 12 models
$=$ 4{,}128 judgements---not only the items strict scoring marks wrong. Full coverage is what makes
a strict \emph{false-positive} rate estimable for the first time: the original audit design, by
construction, could only characterise false negatives. CON is excluded from both judges, matching
the pipeline's existing rationale (Section~\ref{sec:scoring}): it is scored by exact Yes/No match
against an unambiguous binary label, which does not benefit from semantic review. Excluding it is
conservative for the disagreement rates we report below, not convenient: CON is the one task where
strict and judge-based scoring would agree almost by construction, so its absence removes a source
of trivial agreement rather than manufacturing disagreement elsewhere.

Before committing to the full pass we validated the judge on two small, targeted samples. On eight
items Gemini 2.5 Flash itself answered correctly under strict scoring, phi4-mini agreed on all
eight---a sanity check, not evidence of discrimination. On eight items DeepSeek-R1-Distill answered
incorrectly under strict scoring, phi4-mini flipped seven to correct---in every case a genuine
format mismatch (a digit where the reference used a written-out number, or a short preamble before
the answer)---and correctly held the eighth as incorrect, where the model's predicted figure
(``ten per cent'') simply did not match the reference (``twenty per cent''). The judge discriminates
rather than rubber-stamping either the strict score or a blanket verdict.

\textbf{[[JUDGE\_RESULTS\_PENDING]]} Table~\ref{tab:judgereliability} reports, for every model: the
strict false-negative rate (items strict scoring marked wrong that the cross-model judge marks
correct), the strict false-positive rate (the reverse), and agreement between the two judges
(Cohen's $\kappa$). To resolve genuine disagreement between judges from cases where one judge is
simply mistaken, we hand-adjudicate a sample stratified across task type, strict-correct/incorrect
status, and model strength, deliberately oversampling items where the two judges disagree---these
are the only cases informative about which judge is right. Appendix~\ref{app:judge} reports full
per-item results and the manual-audit protocol.

\textbf{Decision rule, stated before the results were examined.} If the cross-model judge supports
DeepSeek-R1-Distill's strict-11th-to-corrected-1st reversal (Section~\ref{sec:regime}), we report it
as the paper's central finding. If it does not, we report the discrepancy plainly and the paper's
central finding becomes the effective-size result (Section~\ref{sec:effsize}) instead. We commit to
this rule here, before running the full pass, specifically so that the choice of headline result is
not made after seeing which one looks better.

\subsection{Statistical Significance and Performance Tiers}
\label{sec:significance}

Paired bootstrap significance testing (10,000 resamples) across all 66 model pairs finds 36 pairs
significantly different at $p < 0.05$, of which 15 survive Bonferroni correction
($\alpha = 0.05/66 \approx 0.00076$). Three descriptive performance clusters organize the results
under strict scoring; their boundaries are robust at the top and bottom of the ranking but
statistical, not categorical, in the middle.

\textbf{Tier 1 --- strong performers (81.5--89.7\%):} Gemini 2.5 Flash, Qwen3-32B, LLaMA-3.3-70B,
Llama 4 Scout 17B, and Kimi K2. \textbf{Tier 2 --- middle performers (75--79\%):} LLaMA-3-8B,
GPT-OSS 120B, GPT-OSS 20B, Gemini 2.5 Pro, Mistral-7B, and DeepSeek-R1-Distill.
\textbf{Tier 3 --- weakest performer:} Gemma 4 E4B, alone at 70.4\%\footnote{[[JUDGE\_RESULTS\_PENDING:
re-verify this tier assignment against the final strict table once the regime section is rebuilt;
Gemma's corrected identity changes its position relative to several Tier 2 models.]]}.

Every Tier 1 model separates from Gemma 4 E4B at the Bonferroni threshold. Within Tier 1, the four
models below Gemini 2.5 Flash are statistically indistinguishable from one another ($p$ = 0.06--0.79),
forming a genuine plateau rather than a fine-grained ranking. Section~\ref{sec:discussion} returns to
two comparisons inside this plateau---Llama 4 Scout 17B vs.\ LLaMA-3.3-70B and GPT-OSS 120B vs.\ 20B---
where a large parameter difference produces no measurable accuracy difference.

\textbf{Robustness to document-level clustering.} 406 items are not 406 independent draws: they trace
to 34 source documents (Section~\ref{sec:models}, Appendix~\ref{app:provenance}), and several
contradiction-detection items link two documents at once. We re-ran \emph{all 66} pairwise
comparisons under a document-clustered bootstrap---resampling the 27 connected components formed by
unioning documents that co-occur in any contradiction item, rather than resampling items
directly---not only the comparisons this paper singles out as findings, since checking only
hand-picked pairs would risk an unrepresentatively reassuring picture. The clustered test finds
\textbf{29 of 66} pairs significant at $p<0.05$ (item-level: 36) and \textbf{11} surviving Bonferroni
correction (item-level: 15--16, itself boundary-sensitive to Monte Carlo noise across independent
10,000-resample runs). \textbf{Nine pairs flip significance direction} under clustering---all among
comparisons this paper does not otherwise discuss. Every comparison this paper specifically reports
as a finding, including the human-reference comparison in Section~\ref{sec:human}, agrees in
direction between the two tests: the correction changes the overall significance count meaningfully
(36$\to$29) without overturning any conclusion this paper draws. Appendix~\ref{app:clustering}
reports the full 66-pair comparison and the clustering construction.
